% section: 等差数列とその総和の導出

\section{等差数列の和}\label{節：等差数列とその総和の導出}
  \sectionref*{節：階差型}で一次式を足すときに使う.
  初項$a$,公差$d$,末項$l=a+(n-1)d$の\bookterm{arithmetic}{等差数列}の和を$S_n$とする.

  \subsection{和の公式}\label{節：等差数列とその総和の導出：総和の公式}
    \[
      S_n=\frac n2(a+l)=\frac n2\{2a+(n-1)d\}.
    \]
    項数と,初項・末項の平均との積である.

  \subsection{公式の導出}\label{節：等差数列とその総和の導出：公式の導出}
    $d>0$の場合,増分$d,2d,\ldots,(n-1)d$を逆順の列と重ねると,各列の高さが$nd$にそろう.
    \[
      2\sum_{k=1}^{n-1}kd=n(n-1)d
    \]
    に初項$a$を$n$回加えれば,$S_n=na+\dfrac{n(n-1)d}{2}$を得る.
    \BookFigArithmeticSum

    同じ工夫を元の数列に使うと,$d$の符号によらず導ける.
    \[
      \begin{array}{rcl}
        S_n&=&a+(a+d)+(a+2d)\\
           &&\qquad+\cdots+\{a+(n-1)d\},\\
        S_n&=&\{a+(n-1)d\}+\{a+(n-2)d\}\\
           &&\qquad+\cdots+(a+d)+a.
      \end{array}
    \]
    各列の和は$2a+(n-1)d$であり,$n$列あるので$2S_n=n\{2a+(n-1)d\}$となる.
